% Section 4: the second output solved on its own terms. The separation lemma is
% what makes it work; the failed greedy variant is kept because it is the
% obvious thing to try and it does not work.
\section{Solving the tree output: the common ancestor}
\label{sec:common-ancestor}

Dijkstra settles most of the maze. The second output makes that unnecessary.

\subsection{Addresses}

Every cell lies in exactly one leaf of the binary tree, and the sequence of
turns from the root to that leaf is the cell's \keyterm{address}
(Figure~\ref{fig:addresses}). Deciding which side of a division a cell falls on
is a comparison of its coordinates against a rectangle --- no cell of the maze
is read, no neighbour list consulted.

% The structure under discussion: a binary tree whose internal nodes hold the
% branching point of the division that made them, and whose leaves are cells.
% A cell's address is the sequence of turns that reaches its leaf.
\begin{figure}[htbp]
  \centering
  \begin{tikzpicture}[scale=1.0, level distance=13mm,
    level 1/.style={sibling distance=44mm},
    level 2/.style={sibling distance=22mm},
    level 3/.style={sibling distance=12mm}]
    \node[tree node] (root) {branching point $(1,2)$}
      child {node[tree node] {$(0,1)$}
        child {node[tree leaf] {\textsf{LL}}}
        child {node[tree leaf] {\textsf{LR}}}
      }
      child {node[tree node] {$(2,3)$}
        child {node[tree leaf] {\textsf{RL}}}
        child {node[tree node] {$(3,1)$}
          child {node[tree leaf] {\textsf{RRL}}}
          child {node[tree leaf] {\textsf{RRR}}}
        }
      };
    \node[font=\scriptsize, black!60, right=3mm of root] {the root region};
  \end{tikzpicture}
  \caption{The proposed structure. Each internal node stores the branching point
  of the division that created it --- the cell its single door sits on --- and
  each leaf is a region of the maze. A cell's \keyterm{address} is the sequence
  of turns from the root to its leaf, so \textsf{RRL} names one region in three
  characters. Sibling leaves are the two halves a single division separated, and
  are joined in the maze by that division's door alone.}
  \label{fig:addresses}
\end{figure}


\subsection{The lemma that does the work}

\begin{proposition}[Separation]
\label{prop:separation}
Let $R$ be an internal node, $R_\ell$ and $R_r$ its children, and $(u,v)$ its
door with $u\in R_\ell$ and $v\in R_r$. For cells $a,b$ of $R$, the unique route
from $a$ to $b$ crosses $(u,v)$ \emph{if and only if} $a$ and $b$ lie in
different children.
\end{proposition}

\begin{proof}
Proposition~\ref{prop:generation-is-a-tree} established that $(u,v)$ is the only
edge between $R_\ell$ and $R_r$, and that each child is spanned by a tree of its
own.

If $a$ and $b$ lie in different children, a route between them uses an edge
between the children, and $(u,v)$ is the only one.

If both lie in the same child, that child's own spanning tree already contains a
route between them. The maze is a tree, so its route between two cells is
unique and is therefore that one; it stays inside the child.
\end{proof}

% The separation lemma in one picture: the door of a division is on the route
% exactly when that division puts the two endpoints on opposite sides.
\begin{figure}[htbp]
  \centering
  \begin{subfigure}[b]{0.47\textwidth}\centering
    \begin{tikzpicture}[scale=0.72]
      \fill[cellgrey] (0,0) rectangle (8,5);
      \draw[black!55] (0,0) rectangle (8,5);
      \draw[wall segment] (4,0) -- (4,2);
      \draw[wall segment] (4,3) -- (4,5);
      \draw[door] (4,2) -- (4,3);
      \node[node vertex, fill=exitorange!25, draw=exitorange] (a) at (1.6,3.8) {};
      \node[node vertex, fill=exitorange!25, draw=exitorange] (b) at (6.4,1.2) {};
      \node[font=\tiny, above=1pt] at (a.north) {$a$};
      \node[font=\tiny, above=1pt] at (b.north) {$b$};
      \draw[linkblue, line width=1pt, rounded corners=3pt]
        (a) -- (2.6,2.5) -- (4,2.5) -- (5.4,2.5) -- (b);
      \node[font=\tiny, doorgreen, right] at (4.15,3.4) {the door};
    \end{tikzpicture}
    \caption{Opposite sides: the route \emph{must} cross, and the door is the
    only crossing there is. Recurse on both halves.}
    \label{fig:separated}
  \end{subfigure}\hfill
  \begin{subfigure}[b]{0.47\textwidth}\centering
    \begin{tikzpicture}[scale=0.72]
      \fill[cellgrey] (0,0) rectangle (8,5);
      \draw[black!55] (0,0) rectangle (8,5);
      \draw[wall segment] (4,0) -- (4,2);
      \draw[wall segment] (4,3) -- (4,5);
      \draw[door] (4,2) -- (4,3);
      \node[node vertex, fill=exitorange!25, draw=exitorange] (a) at (1.2,3.9) {};
      \node[node vertex, fill=exitorange!25, draw=exitorange] (b) at (2.9,1.1) {};
      \node[font=\tiny, above=1pt] at (a.north) {$a$};
      \node[font=\tiny, above=1pt] at (b.north) {$b$};
      \draw[linkblue, line width=1pt, rounded corners=3pt]
        (a) -- (2.2,3.0) -- (1.7,2.0) -- (b);
      \node[font=\tiny, black!55] at (6,2.5) {never entered};
    \end{tikzpicture}
    \caption{Same side: the route stays inside that half, and the other half is
    never looked at. Recurse once.}
    \label{fig:together}
  \end{subfigure}
  \caption{The separation lemma. Whether a division contributes a door to the
  route is decided by a containment test on its two children --- no cell of the
  maze is examined. The right-hand case is where the saving comes from: half the
  region is discarded in one step.}
  \label{fig:separation}
\end{figure}


\begin{keypoint}
The right-hand case is where the saving lives. When both endpoints fall on one
side, \keyterm{half the region is discarded in a single test} --- and the test
never looks at the maze.
\end{keypoint}

\subsection{The algorithm}

\begin{algorithm}[H]
\caption{The route between two cells, read off the binary tree.}
\label{alg:region-route}
\SetKwFunction{Route}{RouteInRegion}
\SetKwFunction{Holds}{Holds}
\KwIn{a node $R$ of the tree, and two cells $a,b$ inside it.}
\KwOut{the route from $a$ to $b$.}
\BlankLine
\Fn{\Route{$R$, $a$, $b$}}{
  \lIf{$R$ \textup{is a corridor leaf}}{\Return the straight run from $a$ to $b$}
  let $(u,v)$ be the door of $R$, with $u\in R_\ell$ and $v\in R_r$\;
  \lIf{$\Holds{$R_\ell$, $a$}$ \textup{\textbf{and}} $\Holds{$R_\ell$, $b$}$}{%
    \Return \Route{$R_\ell$, $a$, $b$}}
  \lIf{$\Holds{$R_r$, $a$}$ \textup{\textbf{and}} $\Holds{$R_r$, $b$}$}{%
    \Return \Route{$R_r$, $a$, $b$}}
  \tcp{separated, so the door is on the route}
  \lIf{$\Holds{$R_\ell$, $a$}$}{%
    \Return \Route{$R_\ell$, $a$, $u$} $\cdot\,(u,v)\,\cdot$ \Route{$R_r$, $v$, $b$}}
  \Return \Route{$R_r$, $a$, $v$} $\cdot\,(v,u)\,\cdot$ \Route{$R_\ell$, $u$, $b$}\;
}
\end{algorithm}

Correctness is Proposition~\ref{prop:separation} applied at every node, with
corridors as the base case: a corridor is a straight line, so the route inside
it is the difference of coordinates.

The first two tests together are a \keyterm{lowest common ancestor} search. The
recursion descends while both endpoints stay on one side and stops descending
at the first node that separates them --- which is their lowest common ancestor,
and whose door is the first door the route crosses.

\begin{description}[leftmargin=*,style=nextline,itemsep=5pt]
  \item[A corridor costs one visit, not its length.] A leaf returns a whole run
    of cells as a single step of the recursion. This is why the node count of
    Section~\ref{sec:analysis} falls \emph{below} the length of the route it
    computes.
  \item[The graph is never read.] Only the tree is. The neighbour lists of
    Section~\ref{sec:maze-generation} are not consulted once.
\end{description}

\subsection{Why not walk the tree greedily}

The obvious alternative is to walk the tree by local moves: from the current
node, step to whichever of the parent, the children and the sibling has its door
nearest the exit. It is worth recording that this does not work, because it is
the first thing one tries.

\begin{center}
\begin{tabular}{@{}lrrr@{}}
  \toprule
  Size & Reached the exit & Never terminated & Stuck \\
  \midrule
  $16\times16$ & $13/200$ & $187$ & --- \\
  $32\times32$ & $6/200$  & $194$ & --- \\
  $64\times64$ & $2/200$  & $198$ & --- \\
  \bottomrule
\end{tabular}
\end{center}

Instrumenting the $32\times32$ failures finds $199$ stopped at a local minimum
and $0$ cycled (Figure~\ref{fig:local-minimum}). Two reasons.

\begin{description}[leftmargin=*,style=nextline,itemsep=5pt]
  \item[It is greedy best-first search.] Ranking moves by an estimate of the
    remaining distance while ignoring the distance travelled is neither optimal
    nor, without a closed list, complete.
  \item[A node's door is not where the walker is.] An internal node stands for a
    \emph{region}; its parent covers the same ground and more. Ranking a parent
    against a child by one stored cell compares things that are not comparable.
\end{description}

% Why the greedy rule stops: the current node is nearer the exit than its
% parent, its two children and its sibling, and is not the goal.
\begin{figure}[htbp]
  \centering
  \begin{tikzpicture}[scale=1.0, every node/.style={font=\scriptsize}]
    \node[tree node] (parent) at (0,2.2) {parent \quad $h=9$};
    \node[tree node, draw=wallred, very thick, fill=wallred!6]
      (here) at (-1.6,0.9) {here \quad $h=\mathbf{4}$};
    \node[tree node] (sib) at (1.9,0.9) {sibling \quad $h=11$};
    \node[tree node] (l) at (-3.0,-0.5) {child \quad $h=7$};
    \node[tree node] (r) at (-0.3,-0.5) {child \quad $h=6$};
    \draw[tree edge] (parent) -- (here);
    \draw[tree edge] (parent) -- (sib);
    \draw[tree edge] (here) -- (l);
    \draw[tree edge] (here) -- (r);
    \draw[wallred, <->, dashed, line width=0.7pt] (here) to[bend right=18] (sib);
    \node[wallred, font=\tiny] at (0.15,1.55) {the shortcut move};
    \node[align=left, font=\scriptsize, black!70] at (3.6,-0.6)
      {every move available\\raises $h$, and this\\node is not the goal};
    \node[wallred, font=\scriptsize\bfseries] at (-1.6,-1.6) {the walk stops here};
  \end{tikzpicture}
  \caption{A local minimum, which is where the greedy rule ends $199$ runs in
  $200$. The stored branching point of the current node happens to lie nearer
  the exit than those of its parent, its two children and its sibling; no move
  improves $h$, and the exit is elsewhere. Nothing about the maze is wrong ---
  the rule is simply hill-climbing on a landscape with hills in it.}
  \label{fig:local-minimum}
\end{figure}


\begin{keypoint}
The estimate answers ``how far is this node's door from the exit''.
Algorithm~\ref{alg:region-route} asks ``which side holds the exit'' --- a
question the tree answers exactly, with no arithmetic at all.
\end{keypoint}
